\newpage
\section*{Day 6 Exercises}
\section*{Exercise 7 --- Simulated Annealing for the TSP}

The Travelling Salesman Problem (TSP) asks for the shortest cyclic tour
visiting $N$ cities exactly once and returning to the starting city. The
tour cost is
\begin{equation}
    f(S) = \sum_{i=1}^{N} D\bigl[S_i,\, S_{i+1 \bmod N}\bigr],
    \label{eq:tourcost}
\end{equation}
where $S$ is a permutation of cities and $D[i,j]$ is the travel cost from
city $i$ to city $j$. Since the number of candidate tours grows as $O(N!)$,
exact optimisation is infeasible for large $N$, and we instead approximate
the optimum using \emph{simulated annealing} (SA).

% ================================================================
\section*{Part 1a --- Circle Sanity Check}
% ================================================================

We place $N = 20$ points evenly on a unit circle, using the Euclidean
distance as $D[i,j]$. The optimal tour is the regular polygon with
analytically known perimeter $2N\sin(\pi/N)$, which provides a ground-truth
check of the implementation.

SA is a stochastic optimisation algorithm that uses the Metropolis--Hastings
acceptance rule with a decreasing temperature schedule. At each iteration
$k$, a neighbouring tour $Y$ is proposed from the current tour $X$. If the
cost decreases ($\Delta = f(Y) - f(X) \le 0$) the move is accepted;
otherwise it is accepted with probability $\exp(-\Delta / T_k)$. This
controlled acceptance of uphill moves prevents the algorithm from getting
trapped in local optima (slide 11), with exploration narrowing as
$T_k \to 0$ (slide 12):
\begin{verbatim}
delta = route_cost(y, D) - route_cost(x, D)
if delta <= 0 or rng.random() < np.exp(-delta / cooling(k)):
    x, cost_x = y, cost_y
    if cost_x < best_cost:
        best_route, best_cost = x.copy(), cost_x
\end{verbatim}
Proposals are generated by swapping two randomly selected cities on the
route, as suggested in the exercise:
\begin{verbatim}
def swap_proposal(route, rng):
    new_route = route.copy()
    i, j = rng.choice(len(route), size=2, replace=False)
    new_route[i], new_route[j] = new_route[j], new_route[i]
    return new_route
\end{verbatim}
We use the cooling schedule $T_k = 1/\sqrt{1+k}$ and run for
$N_{\mathrm{iter}} = 20{,}000$ iterations, tracking the best tour seen.

\begin{table}[h]
    \centering
    \caption{Circle sanity check, $N = 20$.}
    \label{tab:circle}
    \renewcommand{\arraystretch}{1.1}
    \begin{tabular}{lr}
        \toprule
        & Cost \\
        \midrule
        Initial (random) route            & $28.2882$ \\
        Best route (SA)                   & $6.2574$  \\
        Polygon perimeter $2N\sin(\pi/N)$ & $6.2574$  \\
        \bottomrule
    \end{tabular}
\end{table}

SA recovers the exact optimum to four decimal places, confirming the
implementation is correct. Figure~\ref{fig:circle} illustrates the
improvement: the initial route crosses itself many times, while the final
route is the non-self-intersecting regular polygon.

\begin{figure}[H]
    \centering
    \includegraphics[width=0.8\textwidth]{plots/day6/circle_sanity.png}
    \caption{Initial (left) and final (right) routes for 20 points on a unit
             circle. The final route is the non-crossing regular polygon.}
    \label{fig:circle}
\end{figure}

% ================================================================
\section*{Part 1b --- Random Points}
% ================================================================

Now that we have passed the sanity check, we try a more interesting case and draw $N = 20$ points uniformly in the unit square, again using Euclidean distance as $D[i,j]$, the swap proposal, and $T_k = 1/\sqrt{1+k}$ cooling.
As shown in Table~\ref{tab:random}, the tour cost is reduced by 55\% from
the random initial route. Figure~\ref{fig:random} shows that most of the
improvement occurs within the first ${\sim}2{,}000$ iterations, after which
the trajectory flattens as the temperature cools and fewer uphill moves are
accepted.

\begin{table}[H]
    \centering
    \caption{Random points, $N = 20$.}
    \label{tab:random}
    \renewcommand{\arraystretch}{1.1}
    \begin{tabular}{lr}
        \toprule
        & Cost \\
        \midrule
        Initial (random) route & $9.5276$ \\
        Best route (SA)        & $4.3009$ \\
        \bottomrule
    \end{tabular}
\end{table}

\begin{figure}[H]
    \centering
    \includegraphics[width=0.8\textwidth]{plots/day6/random_route.png}
    \caption{Final route (left) and cost-vs-iteration trajectory (right) for
             20 uniform random points in the unit square.}
    \label{fig:random}
\end{figure}
\newpage
% ================================================================
\section*{Part 2 --- General Cost Matrix}
% ================================================================

We now modify the program to work with general travel costs. We load the
$20 \times 20$ asymmetric cost matrix from Learn and initialise SA from the
identity route $0, 1, \ldots, 19$ with the swap proposal and
$T_k = 1/\sqrt{1+k}$ cooling. Since Equation~\eqref{eq:tourcost} already
uses directional lookups $D[S_i, S_{i+1}]$, no modification is needed to
handle asymmetry. The cost computation is vectorised as:
\begin{verbatim}
def route_cost(route, D):
    return D[route, np.roll(route, -1)].sum()
\end{verbatim}
SA reduces the tour cost from $3{,}404$ to $1{,}339$, with best route
(1-indexed towns):
\[
    13,\ 18,\ 7,\ 17,\ 3,\ 20,\ 11,\ 16,\ 10,\ 9,\
    19,\ 15,\ 4,\ 1,\ 8,\ 5,\ 2,\ 14,\ 12,\ 6.
\]

To experiment with different design choices, we implement a
second cooling schedule $T_k = 1/\ln(2+k)$ and a segment-reversal (2-opt)
move, which reverses a randomly chosen contiguous sub-sequence of the tour:
\begin{verbatim}
def reverse_proposal(route, rng):
    i, j = sorted(rng.choice(len(route), size=2, replace=False))
    new_route = route.copy()
    new_route[i:j+1] = new_route[i:j+1][::-1]
    return new_route
\end{verbatim}
We cross both proposals with both cooling schedules and average the best
cost over 5 independent seeds, as reported in Table~\ref{tab:experiment}.

\begin{table}[h]
    \centering
    \caption{Mean best cost over 5 seeds, $N_{\text{iter}} = 20{,}000$.}
    \label{tab:experiment}
    \renewcommand{\arraystretch}{1.1}
    \begin{tabular}{llr}
        \toprule
        Cooling & Proposal & Mean best cost \\
        \midrule
        $1/\sqrt{1+k}$ & Swap    & $1201.0$ \\
        $1/\sqrt{1+k}$ & Reverse & $811.8$  \\
        $1/\ln(2+k)$   & Swap    & $1201.0$ \\
        $1/\ln(2+k)$   & Reverse & $811.8$  \\
        \bottomrule
    \end{tabular}
\end{table}

The proposal mechanism is the dominant factor: segment-reversal achieves a
32\% lower cost than random swaps regardless of the cooling schedule, while
the two cooling schedules perform identically within the iteration budget
used. This is consistent with the practical considerations on candidate
generation (slide 14): the quality of the neighborhood structure drives
performance more than the precise cooling rate. 

\newpage
\section*{Exercise 8 --- Bootstrapping}


In this exercise we apply the \emph{bootstrap}, a non-parametric technique for
estimating the variance of a statistical estimator when no analytical
expression is available (slides 6--7). The idea is to replace the unknown
distribution $F$ by the empirical distribution $F_e$ built from the observed
data, simulate many resampled datasets, and read off the variability of the
resulting estimates.

% ================================================================
\section*{Part 1 --- Ross Exercise 13}
% ================================================================

For i.i.d.\ $X_1,\dots,X_n$ with unknown mean $\mu$ and constants $a<b$ we
want to estimate
\begin{equation}
    p = P\Bigl\{a < \tfrac{1}{n}\textstyle\sum_{i=1}^{n} X_i - \mu < b\Bigr\}.
    \label{eq:p13}
\end{equation}

\textbf{(a)} The empirical distribution $F_e$ assigns mass $1/n$ to each
observation (slide 8),
\begin{equation}
    F_e(x) = \frac{1}{n}\sum_{i=1}^{n} \mathbf{1}\{x_i \le x\}.
    \label{eq:edf}
\end{equation}
Since $nF_e(x)\sim\mathrm{Binomial}(n,F(x))$ we have
$\mathbb{E}[F_e(x)] = F(x)$ and
$\mathbb{V}[F_e(x)] = F(x)(1-F(x))/n$, so $F_e$ is an unbiased and
consistent surrogate for $F$ (slide 10). Drawing from $F_e$ amounts to sampling
with replacement from the data (Ross, p.\ 147).

Under $F_e$ the population mean is the sample mean $\bar x$, so $\bar x$ plays
the role of the unknown $\mu$. The bootstrap estimate of $p$ is then: draw $k$
bootstrap samples of size $n$ with replacement, compute each bootstrap mean
$\bar X^*_i$, and estimate $p$ by the proportion falling within $(a,b)$ of
$\bar x$:
\begin{equation}
    \hat p = \frac{1}{k}\sum_{i=1}^{k}
             \mathbf{1}\bigl\{a < \bar X^*_i - \bar x < b\bigr\}.
    \label{eq:phat}
\end{equation}
The resampling is implemented as:
\begin{verbatim}
def bootstrap_replicates(data, statistic, k, rng):
    n = len(data)
    samples = rng.choice(data, size=(k, n), replace=True)
    return np.apply_along_axis(statistic, 1, samples)
\end{verbatim}

\textbf{(b)} We estimate p, with bootstrap samples $n=10$, data $56,101,78,67,93,87,64,72,80,69$, $a=-5$,
$b=5$ and $k=10{,}000$:

\begin{table}[H]
    \centering
    \caption{Bootstrap estimate of $p$ in Equation~\eqref{eq:p13}.}
    \label{tab:part1}
    \renewcommand{\arraystretch}{1.1}
    \begin{tabular}{lr}
        \toprule
        Quantity & Value \\
        \midrule
        $\bar x$  & $76.70$ \\
        $\hat p$  & $0.7658$ \\
        \bottomrule
    \end{tabular}
\end{table}

So the bootstrap probability that the sample mean of $10$ observations lands
within $\pm 5$ of the true mean is approximately $0.77$.
Figure~\ref{fig:part1} shows the bootstrap distribution of
$\bar X^* - \bar x$ with the interval $[a,b]$ shaded.

\begin{figure}[H]
    \centering
    \includegraphics[width=0.6\textwidth]{plots/day6/ex13_bootstrap_mean.png}
    \caption{Bootstrap distribution of $\bar X^* - \bar x$ ($k=10{,}000$). The
             shaded band is $[a,b]=[-5,5]$; its mass is $\hat p = 0.7658$.}
    \label{fig:part1}
\end{figure}

% ================================================================
\section*{Part 2 --- Ross Exercise 15}
% ================================================================

We want to estimate $\mathrm{Var}(S^2)$, the variance of the sample variance
$S^2 = \sum_{i=1}^{n}(X_i - \bar X)^2/(n-1)$, for $n=15$ and data
$5,4,9,6,21,17,11,20,7,10,21,15,13,16,8$. No convenient closed form exists, so
we apply the bootstrap (slide 11): resample with replacement, recompute
$S^{2*}$ for each bootstrap sample, and take the sample variance of the
$k=10{,}000$ replicates:
\begin{equation}
    \bar\theta^* = \frac{1}{k}\sum_{i=1}^{k}\hat\theta^*_i,
    \qquad
    \hat{\mathbb{V}}[\hat\theta] =
        \frac{1}{k-1}\sum_{i=1}^{k}
        \bigl(\hat\theta^*_i - \bar\theta^*\bigr)^2.
    \label{eq:boot-var}
\end{equation}
This is implemented as seen below and our results are reported in table~\ref{tab:part2}:
\begin{verbatim}
def bootstrap_variance(data, statistic, k, rng):
    replicates = bootstrap_replicates(data, statistic, k, rng)
    return statistic(data), replicates.var(ddof=1)
\end{verbatim}

\begin{table}[H]
    \centering
    \caption{Bootstrap estimate of $\mathrm{Var}(S^2)$.}
    \label{tab:part2}
    \renewcommand{\arraystretch}{1.1}
    \begin{tabular}{lr}
        \toprule
        Quantity & Value \\
        \midrule
        Observed $S^2$ & $34.31$ \\
        $\widehat{\mathrm{Var}}(S^2)$ & $57.92$ \\
        \bottomrule
    \end{tabular}
\end{table}


% ================================================================
\section*{Part 3 --- Sample mean vs.\ sample median}
% ================================================================

We write a program that takes a sample as input and returns the sample median
together with a bootstrap estimate of its variance based on $k=100$ replicates
(Equation~\eqref{eq:boot-var}). We then run it on $n=200$ observations from a
Pareto($\beta=1, k=1.05$) distribution, sampled by inversion
$X = \beta\, U^{-1/k}$ (as on Day 2).

\textbf{1--3)} We compute the sample mean and median, and bootstrap the
variance of each. Our results are reported in table ~\ref{tab:part3}:

\begin{table}[H]
    \centering
    \caption{Bootstrap variance of the mean and the median,
             Pareto($\beta=1, k=1.05$), $n=200$, $k=100$.}
    \label{tab:part3}
    \renewcommand{\arraystretch}{1.1}
    \begin{tabular}{lrrr}
        \toprule
        Estimator & Estimate & Bootstrap variance & Bootstrap std.\ error \\
        \midrule
        Sample mean   & $7.39$ & $8.642$ & $2.940$ \\
        Sample median & $2.06$ & $0.018$ & $0.136$ \\
        \bottomrule
    \end{tabular}
\end{table}

\textbf{4)} The sample median is far more precise  than the sample
mean for this distribution, with a bootstrap variance roughly $8.642 / 0.018 \approx 480$ times
smaller.

\textbf{5)} The difference can be explained by the underlying distribution. For
$k=1.05$ the Pareto mean is finite ($\beta k/(k-1)=21$) but the variance is
\emph{infinite} ($k<2$), so the distribution is extremely heavy-tailed. A few
rare, very large observations dominate the sample mean, making it both imprecise
and biased low for this sample size ($7.39$ versus the true mean of $21$). The
median depends only on the central order statistics and is unaffected by extreme
tail values, giving a much smaller and more stable bootstrap variance. This is a good example of when the bootstrap is useful (slide 5), since
there is no simple formula for the variance of the median.